\documentclass[11pt,a4paper]{article}
\usepackage{amsmath,amssymb,amsthm,mathtools,booktabs,microtype}
\usepackage[margin=25mm]{geometry}
\usepackage[colorlinks=true,linkcolor=blue,urlcolor=blue]{hyperref}
\setlength{\emergencystretch}{2em}
\newtheorem{theorem}{Theorem}[section]
\newtheorem{lemma}[theorem]{Lemma}
\newtheorem{proposition}[theorem]{Proposition}
\newcommand{\Q}{\mathbb Q}\newcommand{\Z}{\mathbb Z}
\DeclareMathOperator{\Tr}{Tr}\DeclareMathOperator{\Nm}{Nm}
\title{Paired contact coordinates:\\an ordinary-abc quartic in 56 variables}
\author{\(\Z\Q\) Diophantine Project}
\date{19 November 2026}
\begin{document}\maketitle
\begin{abstract}
A nonlinear change of contact coordinates reduces the audited
58-variable ordinary-abc quartic to 56 variables. Eight shifts paired
about a fixed center turn the sixteen linear contact factors into four
linear forms in two polynomial invariants. Dividing their product by the
nonzero midpoint allows the original derivative quotient to be encoded
without a further product coordinate. Equal sums of squared radii make
the two paired derivatives differ by a linear form, so only one derivative
coordinate is stored. The resulting circuit has nine contact variables instead of eleven.
A two-adic integrality lemma for the scaled elliptic trace permits a
fixed factor-two padding of the height parameter, preserving the exact
pole square class needed for integer completeness. The polynomial is
a sum of 46 squares of integral quadrics, has 809 coordinate monomials,
degree four, and content one. Under ordinary rational abc and the
explicitly retained arithmetic premises, every sufficiently large odd
shift defines the integers. The sufficient shift remains ineffective.
\end{abstract}

\section{Statement, attribution, and retained inputs}

The shared detector, translated-secant chart, reference rechart, and
Laurent packing used below originate in the notes of \textbf{Gabe Conant
with Astra 6.2}, supplied as \texttt{H10-quartic-constructions.zip} and
\texttt{H10-eleven-variable-constructions.zip}. Their 62-coordinate
construction was audited and extended in v78~\cite{Conant,v78}.
The derivative transfer was proved in v79~\cite{v79}, and its one-helper
normalization and exact pole input were proved in v80~\cite{v80}. The
nonlinear paired-contact rechart here is subsequent project work. The
pole-completeness argument below uses the stated square classes in
Burungale~\cite[Theorem B]{Burungale}.

We retain the parent assertions concerning the elliptic group over
$K=\Q(\sqrt2)$, the inverse-parameter extraction theorem and its formal
comparison estimates, the trace-index bound for all allowed points, the
simultaneous prime-pattern theorem with finite avoidance and local and
real preparation. Ordinary rational abc supplies both the derivative
contact estimate and the signed mixed-power estimate. These retained
premises are not claimed to follow from abc. For the native pole we use
Burungale's Theorem B with its published congruences~\cite{Burungale};
no assertion for the square class $1$ replaces its required class $7$.

\begin{theorem}\label{thm:main}
Under these retained premises and ordinary rational abc, there is an
ineffective constant $B_0$ such that, for every odd integer $B\ge B_0$,
the polynomial specified below has rational input projection $\Z$.
It has 56 total variables, including the input $X$, and is the sum of
46 squares of integral quadrics. Its degree is exactly four, its content
is one, and it has 809 coordinate monomials. For every positive integer
$B$, its exact coefficient height is
\[
 H(F_B)=B^2+377888002299.
\]
This coefficient formula is algebraic and is not a sufficient defining
threshold for $B$.
\end{theorem}

The complete expanded coefficient list and quadratic specification are
\nolinkurl{computation/variable_reductions_v81/abc/quartic56/polynomial.json}
and \nolinkurl{quadrics.json} in that directory. A finite reconstruction
does not certify the retained arithmetic premises.

\section{The unchanged elliptic chart}

Put $\rho^2=2$, $\Delta=11-6\rho$, and
\[
 E:y^2=x^3-\Delta x,\qquad R_0=(1+\rho,2\rho).
\]
We use exactly the complete paired chart of v65 as recharted in
v78. For $X\ne0$ it gives finite points $U,V,H\in E(K)$ with the
specified nonzero abscissas and an allowed reference, with $U\ne V$.
It includes tangencies, zero ordinates, and every finite two-torsion
difference. Its physical outputs are
\[
\begin{gathered}
 s=X\,y(U)/x(U),\qquad t=X\,y(V)/x(V),\\
 S_0=y(H)/x(H),\qquad T_0=y(H-R_0)/x(H-R_0),\\
 J=X^2f(H)f(U),\qquad K'=X^2f(H-R_0)f(V),\qquad f=1/x.
\end{gathered}
\]
Let $C=s-S_0$ and $I=t-T_0$. The stored rational half-trace $h_0$ satisfies
\[
 T_{\mathrm{sc}}=2h_0=X^2\Tr x(U-V).
\]
The retained index $m\in\Z$ of $U-V$ has the parent trace bound
\begin{equation}\label{eq:index}
 m^2\le C_0\bigl(h(T_{\mathrm{sc}})+2h(X)\bigr)+J_0,
\end{equation}
with fixed constants. Finite torsion differences are included with
$m=0$. The parent inverse-parameter extraction theorem says that
positivity of the four valuations of $C,J,\overline I,\overline{K'}$
at one good slot implies that $X$ is locally integral and
\begin{equation}\label{eq:local-label}
                         X(m-X)\equiv0\pmod p.
\end{equation}
All fixed exceptions in these statements remain excluded.

For a nonzero rational $c$ define
\begin{equation}\label{eq:detector}
 W=1+cC^4+c^2J^2-c^3\overline I^{\,4}-c^4\overline{K'}^{\,2},
 \qquad D\Tr W=\Nm(W-1).
\end{equation}
Every rational solution has $\Tr W\ne0$: otherwise the norm vanishes,
so $W=1$ in the field $K$, contrary to $\Tr 1=2$.
The even valuations of reciprocal abscissas give even valuations of
$J,K'$ at every good unramified slot, with zero values interpreted as
having infinite valuation.

\section{Paired contacts and their automatic guards}

Set
\[
 M=B+2\bigl(X^8+T_{\mathrm{sc}}^2\bigr),\qquad
 a=M^4+16,
\]
where $B$ is a fixed positive odd integer. The one free contact parameter
is $z\in\Q$. For the derivative-height theorem introduce the auxiliary,
unrecorded values $x=M^4-z$, $y=M^4+z$, and the fixed shift set
\[
 \mathcal C=\{8,9,12,15,17,20,23,24\},\qquad
 G(t)=\prod_{r\in\mathcal C}(t+r).
\]
All eight shifts are distinct, positive, and have no opposite pair.
Their pairing about 16 gives
\begin{equation}\label{eq:contact-factorization}
 \begin{split}
 F(a,z)&=G(x)G(y)\\
 &=\prod_{r\in\{1,4,7,8\}}\bigl((a-z)^2-r^2\bigr)
                    \bigl((a+z)^2-r^2\bigr).
 \end{split}
\end{equation}
Differentiation in $a$ is the same as differentiation in the original
midpoint $M^4$, with $z$ fixed. Define
\begin{equation}\label{eq:derivative}
 F_a=\frac{\partial F}{\partial a},\qquad
                 u=-16F/F_a,\qquad c=1/u.
\end{equation}

The nonlinear invariants
\[
 T=a^2-z^2,\qquad S=a^2+z^2,\qquad P=T^2
\]
turn the product into
\[
 \Phi(P,S)=\prod_{r\in\{1,4,7,8\}}(P-2r^2S+r^4),\qquad F(a,z)=\Phi(P,S).
\]
By the chain rule,
\begin{equation}\label{eq:chain-rule}
 F_a=2a\bigl(2T\Phi_P+\Phi_S\bigr).
\end{equation}
Store the quotient $\overline F=F/a$; the denominator is nonzero because
$a=M^4+16>0$. The contact quotient is then encoded by
\begin{equation}\label{eq:new-quotient}
             2T\Phi_P+\Phi_S+8c\overline F=0.
\end{equation}
Section~\ref{sec:contact-circuit} rewrites this last equation as a
quadratic residual using a single combined derivative coordinate. Multiplying it by $2a$ recovers $F_a+16cF=0$ exactly.

Together with the literal unit equation $cc^{-1}=1$, these equations
force $F F_a\ne0$. Indeed $G(x)=G(y)=0$ would imply
$x+y=-(r+s)<0$, whereas $x+y=2M^4>0$. If just one factor is zero,
its derivative is nonzero because $G$ is squarefree and the other factor
is nonzero. Thus $F_a\ne0$, contrary to $F_a+16cF=0$.
If $F_a=0$, that same relation first gives $F=0$, which was just
excluded. No inverse of $F$, $F_a$, $x$, $y$, $T$, or $z$ is introduced.
In particular the nonlinear coordinates retain their zero fibers.

The intermediate derivative-height theorem of v63~\cite{v63}, before
replacing the larger slot height by half the product height, applies to
these eight positive shifts. For $\epsilon=1/100$ it gives
\[
 \log r^+(u)\ge\frac{197}{100}\max\{h(x),h(y)\}
                         -\frac{201}{100}h(M)-O(1),
\]
where $r^+(u)$ is the product of primes of positive odd valuation outside
a fixed set. This set includes the bad chart primes, shift differences,
nonzero pair sums, and the primes of $16\cdot7$. The constants are
independent of $B$. Since $x+y=2M^4$, we retain
\begin{equation}\label{eq:radical}
             \log r^+(u)\ge\frac{193}{100}h(M)-O(1).
\end{equation}
If $x$ or $y$ is zero, its eight shifted factors are nonzero constants
and the larger-height slot still supplies the simple-prime argument.
The translated center 16 changes only the fixed shift set, not the
height coefficient or the exponent four.

Since $M\ge B$, estimate~\eqref{eq:radical} forces at least one positive
odd valuation for every allowed rational tuple once $B$ is sufficiently
large. In particular $u$ is not a negative rational square. This
consequence is independent of the pole equations.

\section{One helper, both signs of odd valuation}

Define
\begin{equation}\label{eq:normalizer}
 \alpha=u+u^2D^2,\qquad
 \tau=\frac{\alpha}{1+\alpha},\qquad
 b=7c\tau^2=\frac7u\left(\frac{\alpha}{1+\alpha}\right)^2.
\end{equation}
The gate is the polynomial equation $\tau(1+\alpha)=\alpha$.
It excludes $\alpha=-1$ identically. If $\alpha=0$, then
$u=-(uD)^2$, contrary to~\eqref{eq:radical} for sufficiently large $B$.
Thus $\tau$ and $b$ are nonzero on every sufficiently large defining fiber.
No additional nonzero coordinate is needed.

The native pole equations are
\begin{equation}\label{eq:pole}
 e^2+bp^2=1,\qquad u_0v_0=e(1-b),\qquad u_0^2-v_0^2=bq^2.
\end{equation}
Here $u$ is the contact quotient; $u_0,v_0$ are pole coordinates.

\begin{lemma}\label{lem:pole-obstruction}
If a nonzero rational $b$ admits a rational solution of~\eqref{eq:pole},
then $v_p(b)$ is never a negative odd integer, for any prime $p$.
No nonzero hypothesis on $e,p,u_0,v_0,q$ is needed.
\end{lemma}
\begin{proof}
Suppose $v_p(b)<0$ is odd. The two terms on the left of the first
equation have distinct valuation parity. They cannot cancel, so both
are integral. The term $bp^2$ cannot have valuation zero; hence $e$
is a unit. The middle equation then gives
$v_p(u_0)+v_p(v_0)=v_p(b)$, so both factors are nonzero and their
valuations are unequal. Thus $u_0^2-v_0^2$ has even valuation, whereas
$bq^2$ has odd valuation if $q\ne0$ and is zero if $q=0$.
Both alternatives are impossible.
\end{proof}

\begin{lemma}\label{lem:normalizer}
Let $p\ne7$, let $\mu=v_p(u)$ be odd, and put $d=v_p(D)$, allowing
$d=+\infty$. Then
\[
 A:=v_p(\alpha)=\min\{\mu,2\mu+2d\}.
\]
If $\mu<0$, then $v_p(\tau)=0$ and $v_p(b)=-\mu>0$.
If $\mu>0$ and the native pole equations hold, then
$A>\mu/2$ and $d>-3\mu/4$.
\end{lemma}
\begin{proof}
The two summands defining $\alpha$ have different valuation parity;
hence their minimum is attained once. In particular $\alpha\ne0$.
For $\mu<0$ we have $A\le\mu<0$, so $1+\alpha$ has valuation $A$
and $\tau$ is a unit. This proves the first conclusion for arbitrary $D$.

Suppose $\mu>0$. If $A<0$, then $v_p(\tau)=0$ and $v_p(b)=-\mu$,
contradicting Lemma~\ref{lem:pole-obstruction}. If $A=0$, the gate
ensures its denominator is nonzero, with $v_p(1+\alpha)\ge0$.
Thus $v_p(\tau)\le0$, again giving a negative odd valuation of $b$.
Consequently $A>0$, so $v_p(\tau)=A$ and
$v_p(b)=-\mu+2A$ is odd. It must be positive. Therefore
$A>\mu/2$, which forces $2\mu+2d>\mu/2$, as claimed.
\end{proof}

The negative-odd conclusion is uniform in the actual point denominators.
No second round of local preparation is needed after the auxiliary points
are chosen. The prime seven is included in the fixed exceptional set and
is treated explicitly in the completeness preparation.

\section{Soundness}

Take a consumed prime with $\mu=v_p(u)>0$ odd. Then $v_p(c)=-\mu$.
If both slot valuations of $W$ are negative, say $w_1,w_2$, then
\[
 v_p\Nm(W-1)=w_1+w_2,\qquad
 v_p\Tr W\ge\min(w_1,w_2),
\]
so $v_p(D)\le\max(w_1,w_2)$. Otherwise a slot is already integral.
Lemma~\ref{lem:normalizer} therefore gives a slot at which
$v(W)>-3\mu/4$. This includes an inert prime with repeated slot
valuations.

At that slot the four nonconstant terms of~\eqref{eq:detector} have
pairwise distinct classes modulo four, namely $-\mu,-2\mu,-3\mu,0$.
Here the second and fourth powers have the required classes because
$v(J),v(K')$ are even. Every negative minimum is unique. It follows
that each nonconstant term has valuation greater than $-3\mu/4$.
In particular
\[
 v(C)>\mu/16,\quad v(J)>5\mu/8,\quad
 v(\overline I)>9\mu/16,\quad v(\overline{K'})>13\mu/8.
\]
All four tests are positive. Zero tests satisfy the inequalities with
infinite valuation. The retained extraction theorem now gives
\eqref{eq:local-label}.

If $X\ne0$ and $X\ne m$, writing $X=A/Q$ in lowest terms gives
\[
 \log r^+(u)\le\log|A(mQ-A)|
             \le2h(X)+\log(1+|m|).
\]
Apply the signed mixed-power estimate~\cite{v58} with exponent eight
and $\kappa=1001/1000$ to $s=X^8+T_{\mathrm{sc}}^2$. Since
$M=B+2s$ and $s\ge0$, we have $h(s)\le h(M)+\log2$ uniformly
in positive $B$. Indeed, in lowest terms $s=N/D$, reduction of
$(BD+2N)/D$ cancels a factor dividing two. Therefore
\[
 h(X)\le\frac{501}{997}h(M)+O(1).
\]
Combining with~\eqref{eq:radical} leaves the strict positive margin
\[
 \frac{193}{100}-\frac{1002}{997}
                       =\frac{92221}{99700}>0.
\]
It follows that $h(M)$ is bounded by a fixed multiple of
$\log(1+|m|)+1$. Also
\[
 2h(T_{\mathrm{sc}})\le h(s)+8h(X)+\log2.
\]
The trace estimate~\eqref{eq:index} then bounds $m$, and therefore
$h(M)$, independently of $B$. This contradicts $h(M)\ge\log B$
for sufficiently large $B$. Thus every rational input solution has
$X=0$ or $X=m\in\Z$.

\section{Two-adic preparation and complete integer witnesses}

\begin{lemma}\label{lem:trace-two}
For every nonzero integer $n$, the rational scaled trace
\[
                 T_{\mathrm{sc}}=n^2\Tr x([n]R_0)
\]
is in $\Z_2$. The fixed zero-input fiber used below has
$T_{\mathrm{sc}}=-37888/7199\in\Z_2$ as well.
Consequently every integer witness has
\[
 M=B+2(n^8+T_{\mathrm{sc}}^2)\in\Z_2^\times,
 \qquad a=M^4+16\in\Q_2^{\times2}
\]
for every odd integer $B$.
\end{lemma}
\begin{proof}
Work in $K_2=\Q_2(\rho)$, $\rho^2=2$, with $v_2(2)=1$.
Its integer ring is $\mathcal O=\Z_2[\rho]$. Direct duplication gives
\[
 [2]R_0=\left(\frac{57-28\rho}{8},
             \frac{-660+417\rho}{32}\right).
\]
The tangent slope is $3-\rho/4$. Thus the abscissa has valuation $-3$,
the ordinate has valuation $-9/2$, and the formal parameter
$\xi=-x/y$ has valuation $3/2$.

The integral Weierstrass equation has the formal expansions
$x(\xi)=\xi^{-2}+\cdots$ and $y(\xi)=-\xi^{-3}+\cdots$, with integral
remaining coefficients. Its formal logarithm has the form
$\xi+\sum_{j\ge2}a_j\xi^j/j$ with $a_j\in\mathcal O$.
At $v_2(\xi)>1$, every higher term has strictly greater valuation than
$\xi$; the logarithm is therefore an isometry there. Integer
multiplication preserves this neighborhood. For even $n\ne0$,
\[
 v_2\bigl(\xi([n]R_0)\bigr)=v_2(n)+\tfrac12,
 \qquad v_2\bigl(n^2x([n]R_0)\bigr)=-1.
\]
Hence $n^2x([n]R_0)\in\tfrac12\mathcal O$. Since
$\Tr(\tfrac12\mathcal O)\subseteq\Z_2$, the desired trace is integral.

For odd $n$, write $[n]R_0=R_0+Q$, where $Q$ is in the same formal
neighborhood or is the identity. Translation by the integral point
$R_0$ has an integral power series for its abscissa in $\xi(Q)$.
To see this directly, insert the integral Laurent expansions of
$x(Q),y(Q)$ into the secant formula: $x(Q)-x(R_0)$ has leading term
$\xi^{-2}$, so its normalized denominator is a unit power series;
the negative powers cancel in the translated abscissa, leaving
$x(R_0)+\xi\mathcal O[[\xi]]$. This argument does not assume smooth
reduction of $R_0$. Thus $x([n]R_0)$ is integral, proving the odd case.

For the zero input use the fixed reference $H=2R_0$. Its two ratios are
\[
 S_0=\frac{-348+129\rho}{164},\qquad T_0=4-2\rho.
\]
The constant coefficient of $2T_0^2/(S_0-T_0)$ is
$-37888/7199$. This is the rational scaled trace of that fiber and is
integral at two. Finally $M$ is a two-adic unit for odd $B$, and
$M^4+16\equiv1\pmod8$, which is a nonzero two-adic square.
\end{proof}

Fix $X=n\ne0$. Use the parent multiplication family with
$U-V=[n]R_0$, so $T_{\mathrm{sc}}$, $M$ and $a=M^4+16$ are fixed
before contact prime selection. The sixteen rational positions
$\pm(a+r)$, $r\in\mathcal R:=\{-8,-7,-4,-1,1,4,7,8\}$,
are distinct, since $a>16$. Apply the retained simultaneous
prime-pattern theorem with all fixed exceptions, bounded real and local
point neighborhoods, and the same compatible comparison-index refinements
as in the parent construction. These are sixteen simple contacts instead
of the two previous eight-contact lists; the retained theorem allows any
fixed finite list. At each selected contact prime exactly one of the
sixteen linear factors vanishes simply, its opposite factor is a unit,
and therefore $v_p(u)=1$ by the derivative identity.

The parent successful-slot preparation gives $v(C),v(\overline I)\ge3$
and $v(J),v(\overline{K'})\ge4$ there. Thus $W\equiv1$. At the other
slot the prepared reciprocal products are units and the fourth term in
\eqref{eq:detector} has the unique minimum $-4$. The shifted norm and
trace give $v_p(D)>0$, including $D=0$. Thus $v_p(\alpha)=1$,
$v_p(\tau)=1$, and $v_p(b)=1$ at every selected prime. The selected
primes avoid seven.

For this fixed integer witness, enlarge the finite preparation set to
$S_n$ by the numerator and denominator primes of the now-fixed contact
positions, their nonzero pairwise differences, and the nonzero rational
$a$. This does not
enlarge the fixed exceptional set in the height theorem. Choose a fixed
positive integer $D_0$ with $D_0^2(a+r)\in\Z$ for every shift and take
the prepared slope in the form
\[
 z=\frac{N}{D_0^2L^2},\qquad
 p_{r,\pm}=N\pm D_0^2(a+r)L^2.
\]
The retained prime supplier gives the sixteen distinct positive primes
$p_{r,\pm}$ with the required finite avoidance. Exactly,
\begin{equation}\label{eq:prime-factorization}
 F=\frac{\prod_{r\in\mathcal R}p_{r,-}p_{r,+}}
                {D_0^{32}L^{32}}.
\end{equation}
The sign is positive because there are eight quadratic factors.

Every other prime is covered by the following three cases. Outside
$S_n$, if $z$ is integral and the prime is not one of the sixteen selected
primes, its denominator is a unit: otherwise that prime divides $N,L$
and hence all sixteen distinct selected primes, a contradiction.
Formula~\eqref{eq:prime-factorization} then gives $v_p(F)=0$.
The derivative is integral, so $v_p(u)\le0$. If $v_p(z)<0$ outside
$S_n$, the leading terms are $F=z^{16}+O(z^{14})$ and
$F_a=-16az^{14}+O(z^{12})$ with $a$ a unit. Their strict dominant
valuations give $v_p(u)=2v_p(z)$, an even integer. At the finitely many
places in $S_n$, including any numerator or denominator prime of $a$,
the prescribed sufficiently deep slope denominator gives instead
\[
                      v_p(u)=2v_p(z)-v_p(a).
\]
Choose the slope depth so this valuation is negative. It may be odd;
the uniform normalizer covers that case. The coefficient 16 cancels
between the quotient numerator and derivative leading term, including
at two. Thus all positive odd valuations of $u$ occur among the selected
primes already checked.

At any negative odd valuation of $u$ outside seven, including an
uncontrolled odd factor of the derivative denominator,
Lemma~\ref{lem:normalizer} applies without any restriction on the actual
point coordinates and gives positive odd valuation of $b$.
If $p\ne7$ and $v_p(u)$ is even, then $v_p(b)=-v_p(u)+2v_p(\tau)$
is even. The exceptional prime seven is prepared below with positive
valuation of $b$. Consequently the actual $b$ has square reduced
denominator once that fixed preparation is imposed.

At fixed places and real infinity in the slope, the centered factorization gives
\[
 u\sim z^2/a,\qquad c\sim a/z^2.
\]
The already fixed point neighborhoods keep the raw tests bounded.
Consequently $W-1=O(c)$, $D=O(c^2)$, and
\[
 \alpha=c^{-1}+O(c^2),\qquad \tau=1+O(c),\qquad b\sim7c.
\]
At the real place take $z$ sufficiently large; then $c>0$, $\alpha>0$,
$0<\tau<1$, and $0<b<1$. At two, Lemma~\ref{lem:trace-two} gives $a\in\Q_2^{\times2}$.
Prescribing an arbitrarily deep slope denominator makes $c$ a two-adic
square of arbitrarily large positive even valuation: the ratio
$c/(a/z^2)$ is sufficiently close to one. The quantity $\tau$ is a
nonzero unit. Therefore
\[
 b\in7\Q_2^{\times2},\qquad v_2(b)\ge4.
\]
At every other fixed place, the same preparation makes $v_p(c)$ positive
and $\tau$ a unit; no parity assertion for these valuations is needed. In particular $v_7(b)=1+v_7(c)>0$.
These preparations precede selected prime choices and impose only the
permitted finitely many local conditions.

The next lemma therefore applies to the actual detector value $b$.

\begin{lemma}\label{lem:pole-completeness}
Suppose that $0<b<1$, that $b$ has square reduced denominator, and that
$b\in7\Q_2^{\times2}$ with $v_2(b)\ge4$. Then~\eqref{eq:pole} has a
rational solution.
\end{lemma}
\begin{proof}
Write $b=N/Q^2$ in lowest terms with $0<N<Q^2$ and $Q>0$.
The two-adic assumptions make $Q$ odd and $16\mid N$.
Let $\beta,\eta$ be the positive squarefree parts of $N,Q^2-N$,
respectively. Thus
\[
 b=\beta z_0^2,\qquad 1-b=\eta d_0^2
\]
for nonzero rational $z_0,d_0$. We have $\gcd(\beta,\eta)=1$,
$\beta\equiv7\pmod8$, and $\eta\equiv1\pmod8$.
If $p\mid\beta$, then $1-b\equiv1\pmod p$, so $\eta$ is a square
in $\Q_p$. If $p\mid\eta$, then $b\equiv1\pmod p$, so $\beta$
is a square in $\Q_p$. These are odd primes, so these residue statements
lift to the required local squares. The value $\eta=1$ is allowed.

Apply precisely Theorem B of Burungale~\cite{Burungale}. It supplies
an odd positive squarefree $n\equiv5\pmod8$ and points on
$E_n$ and $E_{n\eta}$ whose nonzero abscissas lie in
$\beta\Q^{\times2}$, where $E_l:y^2=x^3-l^2x$.
Their ordinates are nonzero: the positive nonzero two-torsion abscissas
are $n$ and $n\eta$, with two-adic square class $5$, whereas the
required class is $\beta\equiv7\pmod8$. Negative and zero
abscissas are excluded by $\beta\Q^{\times2}$.
Scale the second point by $(x,y)\mapsto(d_0^2x,d_0^3y)$.
We obtain points
\[
 (ba_1^2,y_1)\in E_n(\Q),\qquad
 (ba_2^2,y_2)\in E_{n(1-b)}(\Q),\qquad a_1a_2\ne0.
\]
Define
\[
 e=\frac{n}{ba_1^2},\qquad p=\frac{y_1}{b^2a_1^3},\qquad
 u_0=\frac{a_2}{a_1},\qquad
 v_0=\frac{e(1-b)}{u_0},\qquad
 q=\frac{y_2}{b^2a_1^3u_0^2}.
\]
All denominators are nonzero. The first curve equation gives
$e^2+bp^2=1$; the definition of $v_0$ gives the middle pole equation;
and the second curve equation gives $u_0^2-v_0^2=bq^2$.
This proves all three equations for this exact rational input.
\end{proof}

For $X=0$, retain the checked parent zero fiber: $H=2R_0$,
$s=S_0$, $t=T_0$, $F_{\mathrm{pair}}=G_{\mathrm{pair}}=d_{
\mathrm{pair}}=J=K'=0$, and
$h=t^2/(s-t)$, whose denominator is nonzero because
$\Nm(S_0-T_0)=7199/328$. Its scaled trace fixes $M$ before the contact
choice. Every raw test vanishes, so $W=1$, $D=0$, $\alpha=u$,
$\tau=u/(1+u)$, and $b=7c/(1+c)^2$. The same contact, real and
fixed-place preparations and Lemma~\ref{lem:pole-completeness} supply
a rational zero. No input factor is
multiplied into the quartic.

\section{The nine-coordinate quadratic contact circuit}\label{sec:contact-circuit}

All nine new contact coordinates are unscaled:
\[
                  a,z,T,S,P,A_1,A_2,\overline F,\mathcal H.
\]
With $M_2=M^2$, impose the four quadratic rows
\[
 a=M_2^2+16,\qquad S+T=2a^2,\qquad S-T=2z^2,\qquad P=T^2.
\]
There is no division by $z$, $S$, $T$, or $P$.
Pair the radii as $(1,8)$ and $(4,7)$. The corresponding products are
\[
\begin{aligned}
 A_1&=P^2-130PS+4097P+256S^2-8320S+4096,\\
 A_2&=P^2-130PS+2657P+3136S^2-101920S+614656.
\end{aligned}
\]
Both radius pairs have sum of squares 65. Their products of squares are
64 and 784. Therefore the paired derivative expressions
\[
 H_i=2T(A_i)_P+(A_i)_S
\]
satisfy the exact linear difference identity
\[
 H_2-H_1=720(-4T+8S-130).
\]
Only $\mathcal H=H_1$ needs a coordinate. The remaining three rows are
\begin{align*}
 a\overline F&=A_1A_2,\\
 \mathcal H&=4TP-260TS+8194T-130P+512S-8320,\\
 \mathcal H(A_1+A_2)&+720(-4T+8S-130)A_1+8c\overline F=0.
\end{align*}
Every residual is quadratic. The first six rows recover the exact
nonlinear invariants and pair-products; the next two recover $F/a$
and the exact $H_1$. Since
\[
 2T\Phi_P+\Phi_S=H_1A_2+A_1H_2,
\]
the last row is precisely~\eqref{eq:new-quotient}, and it recovers the
original derivative contact. Conversely every physical contact tuple
has a unique lift to these nine coordinates because $a\ne0$.
The rechart therefore covers all rational fibers, including zeros of
individual intermediate coordinates.

The former two-slot chart used eleven contact coordinates. Two algebraic
changes give the saving. First, storing $F/a$ removes the product
coordinate required by $a(2T\Phi_P+\Phi_S)+8cF=0$. Second, the equal
sums of squared radii make the two paired derivatives differ linearly,
so a second derivative coordinate is unnecessary. These are exact
quadratic identities, not a dimension count or an assertion on a dense
open subset. The separate normalizer remains exactly the reviewed v80
one; the unresolved cubic in the earlier normalizer-scaling proposal is
not used here.

\section{Exact compilation and coefficient verification}

Use the original v78 Laurent-scaled coordinates for the reference,
paired points, hyperbolic test pairs, input powers, $W,D$, and the first
native pole. Store $c,c^{-1}$ with their literal unit equation.
In particular the physical-to-stored conventions are
\[
 D=cD_s,\quad b=cb_s,\quad e=ce_s,\quad
 u_0=cu_s,\quad v_0=cv_s,\quad p=p_s,\quad q=q_s,
\]
\[
 p^2=cP_s,\qquad q^2=cQ_s.
\]
The five unchanged pole rows therefore read
\[
 p_s^2=cP_s,\quad e_s^2+b_sP_s=c^{-2},\quad
 u_sv_s=e_sc^{-1}-b_se_s,
\]
\[
 q_s^2=cQ_s,\qquad u_s^2-v_s^2=b_sQ_s.
\]
These equations decode exactly to~\eqref{eq:pole}; in particular
there is no hidden additional square in the physical input $b$.

The retained normalizer coordinate is $\alpha$; $\tau$ is retained
from the parent pole input. Its three rows are
\begin{equation}\label{eq:quadratic-gate}
 \alpha=c^{-1}+D_s^2,\qquad
 \tau(1+\alpha)=\alpha,\qquad b_s=7\tau^2.
\end{equation}
Since $D=cD_s$ and $u=c^{-1}$, they recover
$\alpha=u+u^2D^2$ and $b=7c\tau^2$ exactly.
The middle row forces $1+\alpha\ne0$ without any additional guard.
At large defining shifts, $\alpha\ne0$ follows separately from the
derivative-height estimate, as proved above.

Store $M,M_2$ unscaled with
\[
 M=B+2X_4^2+8h_0^2,\qquad M_2=M^2,
\]
where $X_4=X^4$ is the retained input-power coordinate. The nine contact
rows of the preceding section are already quadratic in their unscaled
coordinates. No Laurent rescaling of a contact value or derivative is
needed.

\begin{samepage}
The exact census is
\begin{center}
\begin{tabular}{lr}
\toprule
Coordinates&Count\\\midrule
Input, reference and paired chart, four tests, input powers&30\\
Shared $W$ and $D$&3\\
One native pole, its input, and $\tau$&9\\
$c,c^{-1}$&2\\
$M,M_2$&2\\
Paired-contact nonlinear chart&9\\
One-helper normalizer $\alpha$&1\\\midrule
Total&56\\\bottomrule
\end{tabular}
\end{center}
\end{samepage}
The complete list has 46 quadrics. Their sum of squares vanishes over
$\Q$ precisely when every residual vanishes. Every change of coordinates
used above has a recovered nonzero denominator, including the native
zero-coordinate cases.

Direct exact expansion gives 809 coordinate monomials and 812
parameter-separated terms. The maximum coefficient independent of $B$
is $125368156160$. The only mixed profiles are
$B^2+377888002299$, $16B+1$, and $4B+1$; the remaining parameter-dependent
profiles are linear and smaller than $16B+1$ for positive $B$.
The constant profile exceeds the fixed maximum, and
\[
 B^2+377888002299-(16B+1)=(B-8)^2+377888002234>0.
\]
This proves the exact height formula in Theorem~\ref{thm:main}.
The coefficient of $X^4$ is one, proving both exact degree four and
content one. All 56 coordinates occur.

The independent checker reconstructs the retained chart from its pinned
parent, decodes every new physical scaling, derives the paired factorization and chain-rule
identity independently, expands every squared row with independent
SymPy polynomial arithmetic, and checks the census, support, content,
height profiles and exact rational height margin. The universal
valuation and completeness arguments are the proofs above, not an
inference from a finite residue sample.

\section{Scope of the improvement}

This is a reduction from 58 to 56 variables for the ordinary-abc quartic
route. It replaces the contact block and changes the fixed coefficient
in the height parameter; it retains the same elliptic chart, the same
v80 normalizer, and the same native pole square class. The trace-integrality
lemma and the exact new contact identities are proved here. The stated
elliptic, formal-comparison, trace-index and simultaneous prime-pattern
premises remain in force; no full audit of those arithmetic foundations
is claimed. This rechart does not by itself reduce the separate eleven-
variable unbounded-degree construction, and no global minimum is asserted.

\begin{thebibliography}{9}\footnotesize\raggedright
\setlength{\itemsep}{1pt}
\bibitem{Conant} G.~Conant, with Astra 6.2,
supplied notes \texttt{H10-quartic-constructions.zip} and
\texttt{H10-eleven-variable-constructions.zip}, November 2026;
particularly \texttt{abc62/PROOF.md} and \texttt{fewvars/BALANCED.md}.
The original provenance is retained in the repository's v78 package.
\bibitem{v80} \(\Z\Q\) Diophantine Project,
\emph{One-helper pole normalization: an ordinary-abc quartic in 58 coordinates},
\nolinkurl{papers/quartic_existence/one_helper58_20261119_v80.tex}.
\bibitem{v79} \(\Z\Q\) Diophantine Project,
\emph{One derivative contact and one pole: an ordinary-abc quartic in 59 coordinates},
\nolinkurl{papers/quartic_existence/single_derivative59_20261119_v79.tex}.
\bibitem{Burungale} A.~Burungale,
\emph{A Diophantine definition of $\Z$ in $\Q$},
arXiv:2611.16130v1, November 2026, Theorem B,
\url{https://arxiv.org/html/2611.16130v1}.
\bibitem{v78} \(\Z\Q\) Diophantine Project,
\emph{Audit and extensions of the 62-coordinate quartic},
\nolinkurl{papers/quartic_existence/abc_quartic_audit_20261119_v78.tex}.
\bibitem{v65} J.~Brox and contributors,
\emph{A quartic in 68 variables from translated secants},
\nolinkurl{papers/quartic_existence/quartic_abc68_20261117_v65.tex}.
\bibitem{v63} \(\Z\Q\) Diophantine Project,
\emph{An eleven-variable derivative detector},
\nolinkurl{papers/conditional_few_variables/derivative_detector11_20261117_v63.tex}.
\bibitem{v58} \(\Z\Q\) Diophantine Project, signed mixed-power height theorem,
\nolinkurl{papers/elliptic_heights/signed_mixed_power_height_20261117_v58.tex}.
\end{thebibliography}
\end{document}
